\chapter*{سخن پایانی}
\addcontentsline{toc}{chapter}{سخن پایانی}
\markboth{سخن پایانی}{سخن پایانی}

این کتاب با یک جمله شروع شد: «سلام، دنیا!». آن روز یک ویرایشگر خالی بود و
سه خط کد. امروز می‌دانید یک برنامه چگونه مقدارها را در حافظه نگه می‌دارد،
چگونه حساب می‌کند و تصمیم می‌گیرد، چگونه یک کار را بی‌خستگی تکرار می‌کند و
چگونه کار بزرگ را میان کارکردهای کوچک و نام‌دار تقسیم می‌کند. با آرایه‌ها
مجموعه‌ای از داده‌ها را و با رشته‌ها واژه‌ها و جمله‌ها را در دست گرفته‌اید و
در فصل آخر همه‌ی این‌ها را کنار هم گذاشتید تا برنامه‌هایی بسازید که به
کار می‌آیند.

مهم‌تر از هر ابزار، شیوه‌ی نگاه است که آموختید: مسئله‌ی بزرگ را به قدم‌های
کوچک بشکنید، هر قدم را جداگانه آزمایش کنید، و وقتی برنامه درست کار نکرد،
به پیام خطا به چشم یک راهنما نگاه کنید، نه یک ایراد.

\section*{این کتاب همچنان کنار شماست}

سه پیوست پایانی برای همین ساخته شده‌اند که پس از خواندن کتاب هم به کارتان
بیایند. هر وقت مطمئن نبودید یک واژه رزرو شده است یا نه، به پیوست
\ref{app:keywords} سر بزنید. هر وقت برنامه‌تان پیامی داد که نمی‌فهمیدید،
پیوست \ref{app:errors} را باز کنید. و هر وقت تمرینی را نوشتید و خواستید
راه‌حل خود را با راه‌حلی دیگر بسنجید، پاسخ‌ها در پیوست \ref{app:solutions}
هستند. یادتان باشد که یک تمرین بیش از یک پاسخ درست دارد.

و اگر می‌خواهید همین حالا چیزی بنویسید، ویرایشگر برخط سلام همیشه باز است:
\begin{center}
\url{https://ediotr.salamlang.ir}
\end{center}

\Needspace{18\baselineskip}
\section*{آخرین برنامه}

هر سفری را باید با یک برنامه‌ی سلام بست. این هم آخرین برنامه‌ی کتاب، که
همان برنامه‌ی نخست است با یک خط بیشتر:

\begin{salamfa}
کارکرد آغازین:
    چاپ "سلام، دنیا!"
    چاپ "تا دیدار، برنامه‌نویس!"
پایان
\end{salamfa}

\begin{outputfa}
سلام، دنیا!
تا دیدار، برنامه‌نویس!
\end{outputfa}

هر برنامه‌نویسی روزی با همین «سلام، دنیا!» شروع کرده است. شما هم شروع
کرده‌اید؛ ادامه‌ی راه با شماست.

\vspace{1.5em}
\begin{center}
{\color{brand}\rule{2.5cm}{1pt}}\par\vspace{0.6em}
{\salamcodefont\bfseries\Large\color{salamKeyword}پایان}
\end{center}
